\documentclass[10pt]{article}
\usepackage[a4paper,margin=0.45in]{geometry}
\usepackage{amsmath,amssymb,xcolor}
\setlength{\jot}{4pt}
\allowdisplaybreaks
\begin{document}
\begin{center}{\large\bf Three-$\rho$ terms of the limiting-case cross section ($\Lambda\to0$, $\vee\to\infty$):\\ explicit cancellation and what survives}\end{center}

\section*{0. Result in one paragraph}
In this limit, $\log(1-k^{+}/\vee)\to0$, so every $\log(1-\kappa)$ term is absent. The three-$\rho$ terms fall into two families with identical kernels, plus three rows that stand alone.
\begin{itemize}
\item \textbf{Family A} (kernel $\log u\,\mathcal K_{\rm IV}$): Group I Rows III, IV, Group II Row I (first part), Group III Rows I, II. Every colour structure containing $U(z)$ cancels exactly. The remainder has no $z$ in the colour factor, and it has exactly the colour structure of Group I Row II.
\item \textbf{Family B} (identical $p^{+}$ kernels): Group III Rows III (first part), V, VI. The $U(z)U(y)U(x)$ terms cancel, the $x'\to z$ subtraction terms vanish identically, and the rest collapses to $f^{abc}$ times three LO emission amplitudes.
\item \textbf{The rest} (Group I Row II, the non-$\log u\,\mathcal K_{\rm IV}$ part of Group I Row III, Group I Row V first part) has no partner with the same kernel.
\end{itemize}
So the finite three-$\rho$ terms do \emph{not} cancel identically. What survives is listed in Section 7. All of it is real: $i\times$real plus c.c.\ gives $\sin$, and the one real-coefficient term ($N_{i\pi}$) gives $\cos$. Reality does not require these terms to vanish: a three-$\rho$ term is not an imaginary term. Unlike the $\log(k^{+}/\Lambda)$ three-$\rho$ terms, which must cancel and do (Evolution\_three\_rho file), nothing forces the finite ones to cancel. The four-$\rho$ terms carry only $\log(\vee/\Lambda)$, so they leave no finite three-$\rho$ remainder that could compensate. For a Gaussian (MV) projectile, $\langle\rho\rho\rho\rangle=0$, and this whole sector drops out of the averaged cross section.

\section*{1. Conventions}
\begin{itemize}
\item $\rho$'s are treated as classical (symmetrised): $\{\rho,\rho\}\to2\rho\rho$. Reordering commutators $[\rho^{a}(x),\rho^{b}(y)]=if^{abc}\delta^{(2)}(x-y)\rho^{c}(x)$ produce two-$\rho$ terms with $x=y$; they belong to the two-$\rho$ sector and are not treated here.
\item $U$ is the adjoint Wilson line (real, orthogonal): $U^{cb}(z)U^{c'b}(z)=\delta^{cc'}$.
\item Notation:
\begin{align}
&X^{a}\equiv\big[U^{ab'}(x')-U^{ab'}(w')\big]\rho^{b'}(x'),\qquad D^{d}(x)\equiv U^{de}(x)\rho^{e}(x),\qquad u=\frac{(x-z)^{2}}{(x-w)^{2}},\notag\\
&\mathfrak b^{a}_{i}(w)\equiv\int_{x}\frac{(x-w)_{i}}{(x-w)^{2}}\big[U^{ab}(w)-U^{ab}(x)\big]\rho^{b}(x)\quad\text{(the LO emission amplitude at $w$)},\notag\\
&\mathcal K^{ij}_{\rm IV}=\frac{\delta^{ij}}{2(z-w)^{2}}+\frac{(z-w)^{i}(x-z)^{j}}{(z-w)^{2}(x-z)^{2}}-\frac{(x-w)^{i}(z-w)^{j}}{(z-w)^{2}(x-w)^{2}}-\frac{(x-w)^{i}(x-z)^{j}}{(x-w)^{2}(x-z)^{2}},\notag\\
&\mathbb K\equiv\log\frac{(x-z)^{2}}{(x-w)^{2}}\;\frac{(x'-w')^{i}(y-z)^{j}}{(x'-w')^{2}(y-z)^{2}}\;\mathcal K^{ij}_{\rm IV}.
\end{align}
\item Two integrals over the projectile points reduce to $\mathfrak b$:
\begin{equation}
\int_{x'}\frac{(x'-w')^{i}}{(x'-w')^{2}}X^{a}=-\mathfrak b^{a}_{i}(w'),\qquad\int_{y}\frac{(y-z)^{j}}{(y-z)^{2}}\big[U^{bc}(z)-U^{bc}(y)\big]\rho^{c}(y)=\mathfrak b^{b}_{j}(z).
\end{equation}
\end{itemize}

\section*{2. Which terms have three $\rho$'s}
\begin{tabular}{lll}
\textbf{three $\rho$} & Group I: Rows II, III, IV, V (first part) & Group II: Row I (first part)\\
 & Group III: Rows I, II, III (first part), V, VI & \\[2pt]
\textbf{two $\rho$} & Group I: Row V (second part) & Group II: Row I (second part), Row IV\\
 & Group III: Row III (second part) &
\end{tabular}

\section*{3. Family A: the common kernel $\mathbb K$}
Five rows have exactly the kernel $\mathbb K$ and the phase $e^{-ik\cdot(w'-w)}$. In Group I Row III this is the term $2\log\frac{(x-z)^{2}}{(x-w)^{2}}\mathcal K_{\rm IV}$ of the simplified row. In units of $\frac{1}{(2\pi)^{3}}\frac{g^{4}}{k^{+}}\int_{x,y,z,w,x',w'}e^{-ik\cdot(w'-w)}\,\mathbb K$, they are (each $+$c.c.; for Group II, $+N_{1b}=N_{1a}^{\dagger}$):
\begin{align}
\text{GI R III:}&\quad+\frac{i}{8\pi^{5}}\,f^{abd}X^{a}\,U^{de}(x)\big[U^{bc}(z)-U^{bc}(y)\big]\rho^{e}(x)\rho^{c}(y),\notag\\
\text{GI R IV:}&\quad+\frac{i}{8\pi^{5}}\,X^{a}\big[f^{cdc'}U^{ac}(w)U^{bd}(z)U^{be}(y)\rho^{e}(y)\rho^{c'}(x)-f^{abc}U^{be}(y)U^{cd}(x)\rho^{e}(y)\rho^{d}(x)\big],\notag\\
\text{GII R I.1:}&\quad-\frac{i}{16\pi^{5}}\big[U^{ac'}(w')\{\rho^{c'}(x'),\rho^{b}(y)\}-U^{c'b}(z)U^{ad'}(x')U^{c'e'}(y)\{\rho^{d'}(x'),\rho^{e'}(y)\}\big]B^{ab},\notag\\
\text{GIII R I:}&\quad-\frac{i}{8\pi^{5}}\big[U^{ab'}(x')U^{c'd'}(y)U^{c'b}(z)\rho^{b'}(x')\rho^{d'}(y)-U^{d'b}(z)U^{d'e'}(y)U^{ac'}(w')\rho^{c'}(x')\rho^{e'}(y)\big]B^{ab},\notag\\
\text{GIII R II:}&\quad-\frac{i}{8\pi^{5}}\big[U^{c'd'}(y)U^{ab'}(x')U^{c'b}(z)\rho^{d'}(y)\rho^{b'}(x')-U^{ab'}(x')\rho^{b}(y)\rho^{b'}(x')\big]B^{ab},
\end{align}
with $B^{ab}\equiv f^{cbd}U^{ac}(w)\rho^{d}(x)-f^{acd}U^{cb}(z)U^{de}(x)\rho^{e}(x)$.

\subsection*{Step A1: Group II + Group III}
With $\{\rho,\rho\}\to2\rho\rho$, all three rows carry $-\frac{i}{8\pi^{5}}[\ \cdot\ ]B^{ab}$. Name the four colour structures in the first brackets:
\begin{align}
&g_{1}=U^{ac'}(w')\rho^{c'}(x')\rho^{b}(y),\qquad g_{2}=U^{ab'}(x')U^{c'b}(z)U^{c'e}(y)\rho^{b'}(x')\rho^{e}(y),\notag\\
&g_{3}=U^{ac'}(w')U^{d'b}(z)U^{d'e}(y)\rho^{c'}(x')\rho^{e}(y),\qquad g_{4}=U^{ab'}(x')\rho^{b'}(x')\rho^{b}(y).
\end{align}
Then GII $=[g_{1}-g_{2}]$, GIII R I $=[g_{2}-g_{3}]$ and GIII R II $=[g_{2}-g_{4}]$. Adding them, two of the three $g_{2}$ cancel against each other:
\begin{equation}
g_{1}-g_{2}+g_{2}-g_{3}+g_{2}-g_{4}=(g_{2}-g_{3})-(g_{4}-g_{1})=X^{a}\,U^{c'b}(z)U^{c'e}(y)\rho^{e}(y)-X^{a}\rho^{b}(y).
\end{equation}
Hence
\begin{equation}
\text{GII}+\text{GIII R I}+\text{GIII R II}=+\frac{i}{8\pi^{5}}\,X^{a}\big[\rho^{b}(y)-U^{c'b}(z)U^{c'e}(y)\rho^{e}(y)\big]B^{ab}.
\end{equation}

\subsection*{Step A2: expand}
Multiplying out $[\rho(y)-U^{T}(z)U(y)\rho(y)]^{b}B^{ab}$ gives four terms (all $\times X^{a}$):
\begin{align}
(1)&=+f^{cbd}U^{ac}(w)\rho^{b}(y)\rho^{d}(x),\notag\\
(2)&=-f^{acd}U^{cb}(z)U^{de}(x)\rho^{b}(y)\rho^{e}(x),\notag\\
(3)&=-f^{cbd}U^{ac}(w)U^{c'b}(z)U^{c'e}(y)\rho^{e}(y)\rho^{d}(x),\notag\\
(4)&=+f^{acd}U^{cb}(z)U^{c'b}(z)U^{c'e}(y)U^{de'}(x)\rho^{e}(y)\rho^{e'}(x)=+f^{acd}U^{ce}(y)U^{de'}(x)\rho^{e}(y)\rho^{e'}(x).
\end{align}
In (4) we used $U^{cb}(z)U^{c'b}(z)=\delta^{cc'}$, and $U(z)$ disappears.

\subsection*{Step A3: Group I}
Group I Rows III and IV give (all $\times X^{a}$):
\begin{align}
(A)&=+f^{abd}U^{de}(x)U^{bc}(z)\rho^{e}(x)\rho^{c}(y), &
(B)&=-f^{abd}U^{de}(x)U^{bc}(y)\rho^{e}(x)\rho^{c}(y),\notag\\
(C)&=+f^{cdc'}U^{ac}(w)U^{bd}(z)U^{be}(y)\rho^{e}(y)\rho^{c'}(x), &
(D)&=-f^{abc}U^{be}(y)U^{cd}(x)\rho^{e}(y)\rho^{d}(x).
\end{align}

\subsection*{Step A4: cancellations (relabelling dummy indices only)}
\begin{itemize}
\item $(2)$ with $c\to b,\ b\to c$: $\;-f^{abd}U^{bc}(z)U^{de}(x)\rho^{c}(y)\rho^{e}(x)=-(A)$. \hfill $\Rightarrow (A)+(2)=0$.
\item $(3)$ with $b\to d,\ d\to c',\ c'\to b$: $\;-f^{cdc'}U^{ac}(w)U^{bd}(z)U^{be}(y)\rho^{e}(y)\rho^{c'}(x)=-(C)$. \hfill $\Rightarrow (C)+(3)=0$.
\item $(4)$ with $c\to b,\ d\to c,\ e'\to d$: $\;+f^{abc}U^{be}(y)U^{cd}(x)\rho^{e}(y)\rho^{d}(x)=-(D)$. \hfill $\Rightarrow (D)+(4)=0$.
\end{itemize}
Every structure containing $U(z)$ has cancelled. What is left is $(1)+(B)$. Relabelling:
\begin{align}
(1)&=-f^{hbc}U^{ah}(w)\rho^{b}(x)\rho^{c}(y)=-\tfrac12\,f^{hbc}U^{ah}(w)\{\rho^{b}(x),\rho^{c}(y)\},\notag\\
(B)&=+f^{abc}U^{bd}(x)U^{ce}(y)\rho^{d}(x)\rho^{e}(y)=+\tfrac12\,f^{abc}U^{bd}(x)U^{ce}(y)\{\rho^{d}(x),\rho^{e}(y)\},
\end{align}
so
\begin{equation}
(1)+(B)=-\tfrac12\,S^{a},\qquad S^{a}\equiv f^{hbc}U^{ah}(w)\{\rho^{b}(x),\rho^{c}(y)\}-f^{abc}U^{bd}(x)U^{ce}(y)\{\rho^{d}(x),\rho^{e}(y)\},
\end{equation}
which is \emph{exactly} the colour factor of Group I Row II.

\subsection*{Result for family A}
\begin{equation}
\boxed{\begin{aligned}\text{Family A}&=-\frac{1}{(2\pi)^{3}}\,\frac{ig^{4}}{16\pi^{5}}\,\frac{1}{k^{+}}\int_{w,w',x,y,x'}e^{-ik\cdot(w'-w)}\,\frac{(x'-w')^{i}}{(x'-w')^{2}}\,\mathcal J^{i}(x,y,w)\;X^{a}S^{a}+{\rm c.c.},\\
\mathcal J^{i}(x,y,w)&=\int_{z}\log\frac{(x-z)^{2}}{(x-w)^{2}}\,\frac{(y-z)^{j}}{(y-z)^{2}}\,\mathcal K^{ij}_{\rm IV}\end{aligned}}
\end{equation}
The $z$ integral can be done because $S$ does not depend on $z$. It converges: the integrand is $O(1/|z-w|)$ at $z\to w$, $O(1/|z-y|)$ at $z\to y$, $O(\log|x-z|/|x-z|)$ at $z\to x$, and $O(\log|z|/|z|^{3})$ at large $|z|$.

\noindent\textit{Check.} With exact SU(3) adjoint Wilson lines and random $\rho$ at random points, the sum of the five colour factors equals $-\frac12X^{a}S^{a}$ to $10^{-15}$. Individually the five rows are $O(1)$; for example, in one configuration they are $0.31,\,4.54,\,-0.69,\,4.25$ (GI R III, GI R IV, GII, GIII), and their sum is $8.41$.

\section*{4. Family A remainder together with Group I Row II}
Group I Row II has the same colour factor $X^{a}S^{a}$. With $\mathcal R^{i}(x,y,w)\equiv\int d^{2}p\,d^{2}q\,e^{iq\cdot(w-y)}e^{ip\cdot(y-x)}\frac{R^{i}(q,p)}{p^{2}q^{2}}$, the two combine into
\begin{equation}
\boxed{\begin{aligned}\text{Row II}+\text{Family A}&=-\frac{1}{(2\pi)^{3}}\,\frac{g^{4}}{16\pi^{5}}\,\frac{1}{k^{+}}\int_{w,w',x,y,x'}e^{-ik\cdot(w'-w)}\,\frac{(x'-w')^{i}}{(x'-w')^{2}}\\&\qquad\times\Big[i\,\mathcal J^{i}(x,y,w)+\frac{1}{4\pi}\mathcal R^{i}(x,y,w)\Big]X^{a}S^{a}+{\rm c.c.}\end{aligned}}
\end{equation}
Both $i\mathcal J^{i}$ and $\mathcal R^{i}$ are $i\times$real ($R^{i}$ is odd under $(q,p)\to(-q,-p)$), so with the c.c.\ this term is real ($\sin$).

\textbf{They do not cancel each other.} $\mathcal J^{i}$ decays at large $|x-w|$ (at fixed $y-w$): numerically $|\mathcal J|=4.31$, $2.95$, $1.92$, $1.21$, $0.73$, $0.43$, $0.25$ for $|x-w|=1$, $2$, $4$, $8$, $16$, $32$, $64$, i.e.\ like $\log|x-w|/|x-w|$. By contrast, $R^{i}\to-2q^{i}\log(p^{2}/q^{2})$ as $p\to0$, so $\mathcal R^{i}$ is dominated by large $|x-y|$ and does not decay. Their sum is therefore what remains.

\section*{5. Family B: Group III Rows III (first part), V, VI}
These three rows have identical kernels (both sub-terms) and the same overall factor $\frac{ig^{4}}{4\pi^{4}}$ once $\{\rho,\rho\}\to2\rho\rho$ in Row III. Write $F^{c'}=[U^{c'e'}(y')-U^{c'e'}(x')]\rho^{e'}(x')$ (first sub-term) or $F^{c'}=[U^{c'e'}(z)-U^{c'e'}(x')]\rho^{e'}(x')$ (the $x'\to z$ subtraction). The colour factors are:
\begin{align}
\text{R III}&:\ f^{c'd'a}F^{c'}U^{d'b}(z)\big[U^{ac}(w)\rho^{c}(y)\rho^{b}(x)-U^{cb}(z)U^{ad}(y)U^{ce}(x)\rho^{d}(y)\rho^{e}(x)\big],\notag\\
\text{R VI}&:\ f^{c'd'a}F^{c'}U^{d'b}(z)\big[U^{cb}(z)U^{ad}(y)U^{ce}(x)\rho^{d}(y)\rho^{e}(x)-U^{ad}(y)\rho^{d}(y)\rho^{b}(x)\big],\notag\\
\text{R V}&:\ f^{c'd'a}F^{c'}\big[U^{d'e}(x)U^{ad}(y)-U^{ad}(w)U^{d'e}(x)\big]\rho^{e}(x)\rho^{d}(y).
\end{align}
\textbf{Step B1.} In III $+$ VI the $U(z)U(y)U(x)$ terms cancel:
\begin{equation}
\text{III}+\text{VI}=f^{c'd'a}F^{c'}\,U^{d'b}(z)\rho^{b}(x)\,\big[U^{ad}(w)-U^{ad}(y)\big]\rho^{d}(y).
\end{equation}
\textbf{Step B2.} Row V is $-f^{c'd'a}F^{c'}U^{d'e}(x)\rho^{e}(x)[U^{ad}(w)-U^{ad}(y)]\rho^{d}(y)$, so
\begin{equation}
\text{III}+\text{V}+\text{VI}=f^{c'd'a}\,F^{c'}\,\big[U^{d'b}(z)-U^{d'b}(x)\big]\rho^{b}(x)\,\big[U^{ad}(w)-U^{ad}(y)\big]\rho^{d}(y).
\end{equation}
This was checked with exact SU(3) Wilson lines: $25.35178=25.35178$, from $10.41+7.75+7.19$.

\noindent\textbf{Step B3.} The kernel carries $\frac{(x-z)^{j}}{(x-z)^{2}}$, $\frac{(y-w)^{i}}{(y-w)^{2}}$, and $\frac{(x'-y')^{k'}}{(x'-y')^{2}}$ or $\frac{(x'-z)^{k'}}{(x'-z)^{2}}$. So the $x$, $y$, $x'$ integrals give LO amplitudes: $\mathfrak b^{d'}_{j}(z)$, $\mathfrak b^{a}_{i}(w)$, and $\mathfrak b^{c'}_{k'}(y')$ or $\mathfrak b^{c'}_{k'}(z)$. The subtraction sub-term has $j=k'$ (from $\delta_{jk'}$), so its colour is $f^{c'd'a}\mathfrak b^{c'}_{j}(z)\mathfrak b^{d'}_{j}(z)=0$ by antisymmetry. \emph{The $x'\to z$ subtraction terms vanish identically.} Hence
\begin{align}
\boxed{\begin{aligned}
\text{Family B}&=\frac{1}{(2\pi)^{3}}\,\frac{ig^{4}}{4\pi^{4}}\,\frac{1}{k^{+}}\int_{0}^{\infty}\frac{dp^{+}}{2\pi}\int_{w,y',z}e^{-ik\cdot(y'-w)}\,e^{-ik\cdot\frac{p^{+}}{k^{+}}(y'-z)}\,\frac{(y'-z)^{m}}{(y'-z)^{2}}\,f^{c'd'a}\\
&\times\Bigg\{\frac{\mathfrak b^{c'}_{m}(y')\,\mathfrak b^{d'}_{j}(z)\,\mathfrak b^{a}_{j}(w)}{p^{+}+k^{+}}-\Big(\Big[\frac{k^{+}}{p^{+}(p^{+}+k^{+})}\Big]_{+}+\frac{1}{p^{+}+k^{+}}\Big)\mathfrak b^{c'}_{i}(y')\,\mathfrak b^{d'}_{m}(z)\,\mathfrak b^{a}_{i}(w)
-\frac{1}{k^{+}}\,\mathfrak b^{c'}_{j}(y')\,\mathfrak b^{d'}_{j}(z)\,\mathfrak b^{a}_{m}(w)\Bigg\}+{\rm c.c.}
\end{aligned}}
\end{align}
This is the three-gluon vertex $f^{c'd'a}$ contracted with three LO emission amplitudes, at $y'$, $z$ and $w$. It is $i\times$real, so with the c.c.\ it is real.

\section*{6. Terms without a partner}
\textbf{Group I Row III, the part that is not $\log u\,\mathcal K_{\rm IV}$.} Using $D^{d}(x)$ and $\mathfrak b$:
\begin{align}
&-\frac{1}{(2\pi)^{3}}\frac{ig^{4}f^{abd}}{16\pi^{5}}\frac{1}{k^{+}}\int_{w,w',x,z}e^{-ik\cdot(w'-w)}\,\mathfrak b^{a}_{i}(w')\,D^{d}(x)\Bigg\{\bigg[\frac{2\,\theta_{w}}{|(x-w)\times(z-w)|\,(z-w)^{2}}\Big((z-w)^{i}X_{\perp}^{j}+Z_{\perp}^{i}(z-w)^{j}\Big)\notag\\
&\qquad+\frac{1}{(z-w)^{2}}\log\frac{(z-w)^{2}}{(x-w)^{2}}\Big(\frac{(x-w)^{i}(z-w)^{j}}{(x-w)^{2}}-\frac{(z-w)^{i}(x-z)^{j}}{(x-z)^{2}}\Big)\bigg]\mathfrak b^{b}_{j}(z)
-\log\frac{(z-w)^{2}}{(x-w)^{2}}\,\frac{(x-z)\cdot(w-z)}{(x-z)^{2}(w-z)^{2}}\,\mathfrak b^{b}_{i}(w)\Bigg\}\notag\\
&+\frac{1}{(2\pi)^{3}}\frac{g^{4}f^{abd}}{8\pi^{4}}\frac{1}{k^{+}}\int_{w,w',x,z}e^{-ik\cdot(w'-w)}\,\mathfrak b^{a}_{i}(w')\,\mathcal K^{ij}_{\rm IV}\,D^{d}(x)\,\mathfrak b^{b}_{j}(z)\;+\;{\rm c.c.}
\end{align}
The last line is $N_{i\pi}$. It is the only real-coefficient three-$\rho$ term ($\cos$); its $z\to w$ singular part is the Coulomb phase, which cancels against the c.c.

\noindent\textbf{Group I Row II.} It is already combined with family A in Section 4.

\noindent\textbf{Group I Row V, first part.} Unchanged. Its phase $e^{-ik\cdot(w'-z)}e^{-ik\cdot(z-x)p^{+}/k^{+}}$ with $0<p^{+}<k^{+}$ matches no other term.

\section*{7. The complete three-$\rho$ sector}
\begin{equation}
\frac{d^{3}N^{3\rho}}{d^{3}k}=\big[\text{Row II}+\text{Family A}\big]_{\S4}+\big[\text{Family B}\big]_{\S5}+\big[\text{GI R III remainder}\big]_{\S6}+\big[\text{GI R V, first part}\big].
\end{equation}
Compared with the ten three-$\rho$ rows of the file, every structure with $U(z)$ in Family A and every $U(z)U(y)U(x)$ and subtraction structure in Family B has cancelled. Each surviving term is real, by the same mechanism as before. The sector does not cancel to zero.

\section*{8. Two corrections in the uploaded file}
\begin{enumerate}
\item Group I Row II: the first term of $R^{i}$ reads $-2k^{i}\log\frac{p^{2}}{q^{2}}$. It must be $-2q^{i}\log\frac{p^{2}}{q^{2}}$ (the dummy momentum was renamed $k\to q$).
\item Group I Row III, last term of $N_{3b}$: the colour factor is written with $U^{bc}(z)$, but the kernel is the amended one, and your own summed Row III uses $U^{bc}(w)$. It must be $\big[U^{de}(x)U^{bc}(w)-U^{de}(x)U^{bc}(y)\big]\rho^{e}(x)\rho^{c}(y)$, and the denominator should read $(y-w)^{2}$, as in the sum.
\end{enumerate}
\end{document}
